\section*{Chapter \ref{ch:b3:minimum-wage} --- Minimum Wages and Monopsony}

\begin{solution}{exo:b3:minimum-wage:1}
Three of the five. The two New Jersey estimates use the same increase and the same
comparison state; they differ in the data --- a telephone survey of managers against
payroll records. % figdata: elasticities
\end{solution}

\begin{solution}{exo:b3:minimum-wage:2}
$-0.2 \times 10~\% = -2~\%$.
\end{solution}

\begin{solution}{exo:b3:minimum-wage:3}
Marginal cost $4 + 0.4L = 24 - 0.2L$ gives $L = 33.3$ and, on the supply curve, $w = 10.67$. Competitive:
$4 + 0.2L = 24 - 0.2L$, $L = 50$, $w = 14$.
\end{solution}

\begin{solution}{exo:b3:minimum-wage:4}
An employer who loses all its workers when it cuts its wage slightly below other
employers': a perfectly elastic supply to the firm.
\end{solution}

\begin{solution}{exo:b3:minimum-wage:5}
$1/(1 + 4) = 20~\%$ of the marginal revenue product.
\end{solution}

\begin{solution}{exo:b3:minimum-wage:6}
Floor 12: supply gives $(12 - 4)/0.2 = 40$, demand $(24 - 12)/0.2 = 60$: employment 40. Floor 14: 50 on both.
Floor 16: supply 60, demand 40: employment 40. All three exceed the monopsony 33.3.
\end{solution}

\begin{solution}{exo:b3:minimum-wage:7}
The tested solution is \code{ex07_bunching.py}, in the chapter's \code{exercises/}
folder. Net of the control states, 41.0 thousand jobs are missing below 10 and 39.0 thousand
are in excess from 10 to 13: employment falls by 2.0 thousand, 2.1~\% of jobs below 14
dollars. % output: ex07_bunching
\end{solution}

\begin{solution}{exo:b3:minimum-wage:8}
$0.30 \times 10~\% = 3~\%$.
\end{solution}

\begin{solution}{exo:b3:minimum-wage:9}
Imprecise estimates concentrated on one side --- many small studies with large
negative (or large positive) elasticities and few on the other --- while precise
estimates cluster near a different value. Five estimates are too few to show an
asymmetry, and a \omterm{def:b3:credible-empirical-work:meta-analysis}{funnel plot} needs each estimate's standard error, which the
abstracts behind the opening chart mostly do not report.
\end{solution}

\begin{solution}{exo:b3:minimum-wage:10}
The proof is in the text. As $\varepsilon \to \infty$, $\varepsilon/(1 + \varepsilon) \to 1$ and $1/(1 + \varepsilon) \to 0$: the wage equals the marginal
revenue product, as in the competitive model.
\end{solution}

\begin{solution}{exo:b3:minimum-wage:11}
The tested solution is \code{ex11_supply.py}, in the chapter's \code{exercises/}
folder. The slope is 3.97, close to the true 4, and the implied markdown is 20.1~\% of the
marginal revenue product. % output: ex11_supply
Differences in productivity move each firm along the common supply curve, which is
why the regression traces it.
\end{solution}

\begin{solution}{exo:b3:minimum-wage:12}
At a \omterm{def:b3:minimum-wage:kaitz}{Kaitz index} of 0.8 the floor is likely above the competitive wage for many
employers, where both models predict job losses. The most relevant evidence is from
large single increases: Hungary (small losses, prices up), Seattle (hours cut) and
Germany (no losses at a narrower reach); state panels of moderate increases say little.
\end{solution}

\begin{solution}{pb:b3:minimum-wage:brenland}
% figdata: brenland_stats, brenland_hist
\textbf{1.} 16.89 dalens; $12/16.89 = 0.71$.

\textbf{2.} 20.2~\% (290 of 1{,}439).

\textbf{3.} By 34.3~\% on average.

\textbf{4.} $1.0 \times 34.3~\%$ of their jobs, about 100 heads: 6.9~\% of working heads.

\textbf{5.} An elasticity describes small changes; a rise of a third moves far along the
demand curve, where its slope may differ, and many of those affected would leave
for other work rather than unemployment.

\textbf{6.} $6 + 0.2L = 20 - 0.1L$: $L = 46.7$, $w = 6 + 0.1 \times 46.7 = 10.67$.

\textbf{7.} $20 - 0.1 \times 46.7 = 15.33$; the wage falls short by $(15.33 - 10.67)/15.33 = 30.4~\%$.

\textbf{8.} $\varepsilon = 10.67/(0.1 \times 46.7) = 2.29$; $1/(1 + 2.29) = 0.304$, as in question 7.

\textbf{9.} $6 + 0.1L = 20 - 0.1L$: $L = 70$, $w = 13$.

\textbf{10.} Floor 12: $\min(60, 80) = 60$. Floor 15: $\min(90, 50) = 50$. Floor 17: $\min(110, 30) = 30$.

\textbf{11.} The competitive wage, 13: supply and demand both give 70.

\textbf{12.} Above $20 - 0.1 \times 46.7 = 15.33$, where demand falls below the monopsony employment.

\textbf{13.} 290.

\textbf{14.} $0.9 \times 290 = 261$.

\textbf{15.} $261 - 290 = -29$ heads, $-2.0~\%$ of working heads.

\textbf{16.} Jobs pushed into 13--14 would be missed by a window that stops at 13, making
the estimate too negative; the window should include them.

\textbf{17.} Compare the bins above 13 before and after with the same bins in an area
without the floor.

\textbf{18.} $(99.0 - 100) - (100.5 - 100) = -1.5~\%$.

\textbf{19.} $0.343 \times 0.202 = 6.9~\%$, if the wages of others are unchanged.

\textbf{20.} $-1.5/6.9 = -0.22$.

\textbf{21.} Part I predicted a loss of 6.9~\% of working heads; the border comparison finds
1.5~\%, about a fifth.

\textbf{22.} Different trends on the two sides of the border; Kestrian employers raising
wages to keep workers, which contaminates the control.

\textbf{23.} If Brenland's workers take jobs across the border, Kestria's employment rises
and Brenland's falls for reasons other than lost jobs: the estimate is too negative.

\textbf{24.} That they have some power over wages, or adjust through hours and prices
rather than heads.

\textbf{25.} Phase the floor in by region at random dates, and measure heads, hours and
prices against regions not yet covered.

\textbf{26.} In the firm model, the employment-maximising floor is 13 dalens an hour; the
border comparison measures the floor of 12 as lowering employment by 1.5~\%, an own-wage
elasticity of $-0.22$.
\end{solution}
